\documentclass[12pt, a4paper]{article}

% ─── Paketler ───────────────────────────────────────────────────────────────
\usepackage[utf8]{inputenc}
\usepackage[T1]{fontenc}
\usepackage[turkish]{babel}
\usepackage{geometry}
\usepackage{graphicx}
\usepackage{amsmath}
\usepackage{listings}
\usepackage{xcolor}
\usepackage{hyperref}
\usepackage{booktabs}
\usepackage{float}
\usepackage{caption}
\usepackage{subcaption}
\usepackage{fancyhdr}
\usepackage{titlesec}
\usepackage{enumitem}
\usepackage{parskip}

\geometry{
  top=2.5cm, bottom=2.5cm,
  left=3cm, right=2.5cm
}

% ─── Kod renk şeması ────────────────────────────────────────────────────────
\definecolor{codegreen}{rgb}{0,0.6,0}
\definecolor{codegray}{rgb}{0.5,0.5,0.5}
\definecolor{codepurple}{rgb}{0.58,0,0.82}
\definecolor{backcolour}{rgb}{0.95,0.95,0.92}

\lstdefinestyle{cstyle}{
  backgroundcolor=\color{backcolour},
  commentstyle=\color{codegreen},
  keywordstyle=\color{blue},
  numberstyle=\tiny\color{codegray},
  stringstyle=\color{codepurple},
  basicstyle=\ttfamily\footnotesize,
  breakatwhitespace=false,
  breaklines=true,
  captionpos=b,
  keepspaces=true,
  numbers=left,
  numbersep=5pt,
  showspaces=false,
  showstringspaces=false,
  showtabs=false,
  tabsize=2,
  language=C
}

% ─── Üst/alt bilgi ──────────────────────────────────────────────────────────
\pagestyle{fancy}
\fancyhf{}
\rhead{INF 340 Mikroişlemciler}
\lhead{Akıllı Ekosistem -- Son Rapor}
\rfoot{\thepage}

% ─── Başlık ayarları ────────────────────────────────────────────────────────
\titleformat{\section}{\large\bfseries}{}{0em}{\thesection\quad}
\titleformat{\subsection}{\normalsize\bfseries}{}{0em}{\thesubsection\quad}

% ────────────────────────────────────────────────────────────────────────────
\begin{document}

% ─── Kapak ──────────────────────────────────────────────────────────────────
\begin{titlepage}
  \centering
  \vspace*{1.5cm}
  {\Large \textbf{INF 340 MİKROİŞLEMCİLER}\par}
  \vspace{0.5cm}
  {\large Proje Son Raporu\par}
  \vspace{1.5cm}
  {\LARGE \textbf{Akıllı Ekosistem:\\ Otonom İçecek ve Çevre Yönetimi}\par}
  \vspace{2cm}
  \begin{tabular}{ll}
    \textbf{Ayşe Selen Göğüş}  & 23401063 \\[4pt]
    \textbf{Resul Sanberk Kalko} & 23401068 \\[4pt]
    \textbf{Erdem Ertan}         & 22401986 \\
  \end{tabular}
  \vfill
  {\large Mayıs 2026\par}
\end{titlepage}

% ─── İçindekiler ────────────────────────────────────────────────────────────
\tableofcontents
\newpage

% ════════════════════════════════════════════════════════════════════════════
\section{Proje Özeti}
% ════════════════════════════════════════════════════════════════════════════

Bu proje kapsamında, mikrodenetleyici tabanlı bir akıllı ekosistem sistemi
tasarlanmış ve hayata geçirilmiştir. Sistem; termodinamik ve akustik sensörleri bir araya
getirerek sıradan bir yaşam alanını üç temel otonom işlev üzerinden yönetebilen bir
platforma dönüştürmektedir:

\begin{enumerate}[leftmargin=*, label=\arabic*.]
  \item \textbf{Soğuk Algılama ve Müzik Açma:} İçeceğin sıcaklığı hedef değere
        ulaştığında Bluetooth üzerinden otomatik müzik başlatılması.
  \item \textbf{Temassız Jest Kontrolü:} El kaydırma hareketiyle çalan şarkının
        ileri/geri değiştirilmesi.
  \item \textbf{Akustik Lamba Kontrolü:} El çırpma sesiyle ortam lambasının açılıp
        kapatılması ve parlaklığının ayarlanması.
\end{enumerate}

% ════════════════════════════════════════════════════════════════════════════
\section{Sistem Mimarisi}
% ════════════════════════════════════════════════════════════════════════════

Sistemin genel mimarisi Şekil~\ref{fig:mimari} üzerinde gösterilmiştir. MSP432 işlemcisi
merkezi denetleyici görevi üstlenerek tüm sensörlerden veri toplamakta ve ilgili çevresel
birimleri yönetmektedir.

\begin{figure}[H]
  \centering
  % Gerçek sistem diyagramı buraya eklenecek:
  % \includegraphics[width=0.85\textwidth]{mimari.png}
  \fbox{\rule{0pt}{5cm}\rule{0.85\textwidth}{0pt}}
  \caption{Sistem Blok Diyagramı}
  \label{fig:mimari}
\end{figure}

Sistem, sürekli bekleme döngüleri yerine \textbf{kesme (Interrupt) tabanlı} bir mimari
üzerine kurulmuştur. Her sensör olayı (sıcaklık eşiği, optik engelleme, ses darbesi)
ayrı bir donanımsal kesme vektörüne atanmıştır; bu sayede işlemci boşta kalırken
\texttt{\_\_low\_power\_mode\_0()} durumunda bekleyebilmekte ve gereksiz güç tüketimi
önlenmektedir.

\subsection{Durum Şeması}

Sistemin genel akışı aşağıdaki durum geçişleriyle modellenmiştir:

\begin{center}
\begin{tabular}{lll}
  \toprule
  \textbf{Durum} & \textbf{Tetikleyici} & \textbf{Geçilen Durum} \\
  \midrule
  \texttt{IDLE}         & İçecek sıcaklığı < eşik           & \texttt{IDLE} \\
  \texttt{IDLE}         & İçecek sıcaklığı $\geq$ eşik      & \texttt{MUSIC\_ON} \\
  \texttt{MUSIC\_ON}    & Jest algılama (sağ/sol)            & \texttt{MUSIC\_ON} \\
  \texttt{MUSIC\_ON | IDLE }    & El çırpma                          & \texttt{LAMP\_TOGGLE} \\
  \texttt{LAMP\_TOGGLE} & —                                  & \texttt{MUSIC\_ON} \\
  \bottomrule
\end{tabular}
\end{center}

% ════════════════════════════════════════════════════════════════════════════
\section{Kullanılan Donanım}
% ════════════════════════════════════════════════════════════════════════════

\begin{table}[H]
  \centering
  \caption{Donanım Bileşenleri}
  \label{tab:donanim}
  \begin{tabular}{lll}
    \toprule
    \textbf{Bileşen}          & \textbf{Model}      & \textbf{Arayüz} \\
    \midrule
    Ana Kontrolör             & MSP432P401R         & —         \\
    Bluetooth Modülü          & HC-05               & UART      \\
    Sıcaklık Sensörü          & DS18B20             & 1-Wire    \\
    Optik Sensör (jest)       & LDR çifti           & GPIO/ADC  \\
    Mikrofon Modülü           & KY-037              & GPIO (INT)\\
    Röle Modülü               & 5 V Röle            & GPIO      \\
    OLED Ekran                & 0.96" SSD1306       & I2C       \\
    \bottomrule
  \end{tabular}
\end{table}

% ════════════════════════════════════════════════════════════════════════════
\section{Özellik 1: Soğuk Algılama ve Müzik Açma}
% ════════════════════════════════════════════════════════════════════════════

\subsection{Çalışma Prensibi}

\textbf{DS18B20} sıcaklık sensörü, \textbf{1-Wire} protokolü üzerinden \textbf{Arduino Uno}'ya periyodik olarak sıcaklık verisi iletmektedir. Ölçüm döngüsü, ana program akışı içerisinde yazılımsal bir gecikme (\textbf{delay(2000)}) aracılığıyla 2 saniyede bir tetiklenmektedir. Sensörden okunan dijital veriler, harici bir ADC dönüşüm formülüne gerek kalmadan \textbf{DallasTemperature} kütüphanesi kullanılarak doğrudan santigrat dereceye çevrilmektedir.
\\ \\
Dönüştürülen bu değer, sistemde belirlenen hedef sıcaklık eşiği (örneğin 10°C) ile kıyaslanmaktadır. Sıcaklık eşiğinin altına inildiğinde ve ilgili bayrak değişkeni (\textbf{sarkiCalindi} ) kontrol edildiğinde sistem bekleme durumundan çıkarak \textbf{Bluetooth} modülünü tetikler.

\subsection{Bluetooth Haberleşmesi (HC-05 / UART)}

Sistemdeki haberleşme, donanımsal çakışmaları ve zamanlama hatalarını önlemek amacıyla \textbf{SoftwareSerial} kütüphanesi kullanılarak \textbf{D10 (RX)} ve \textbf{D11 (TX) }pinleri üzerinden 9600 baud hızında gerçekleştirilmektedir. Sıcaklık eşiği aşıldığında \textbf{HC-05} modülü üzerinden "\textbf{PLAY\_MUSIC}" komutu gönderilmektedir. Bilgisayar ortamında arka planda çalışan Python uygulaması (\textbf{pyserial} ve \textbf{spotipy} kütüphaneleri ile) bu portu dinlemekte ve komutu yakaladığı an yetkilendirilmiş API üzerinden Spotify medyasını otonom olarak başlatmaktadır.

\subsection{Kritik Kod Parçacığı}

\begin{lstlisting}[style=cstyle, caption={Sıcaklık eşiği kontrolü ve UART tetiklemesi}]
void loop() {
  sensors.requestTemperatures(); 
  float anlikSicaklik = sensors.getTempCByIndex(0);

  // Sicaklik limiti kontrolu ve UART uzerinden Python'u tetikleme
  if (anlikSicaklik <= hedefSicaklik && !sarkiCalindi && anlikSicaklik > -50.0) {
    display.setTextSize(1);
    display.setCursor(0, 50);
    display.println("Soguk! Muzik Caliyor");
    
    btSerial.println("PLAY_MUSIC"); // HC-05 uzerinden komut gonder
    sarkiCalindi = true; 
  } 
  
  display.display();
  delay(2000); 
}
\end{lstlisting}

\subsection{Kalibrasyon ve Test Sonuçları}

Sensör verileri, Arduino IDE Seri Port Ekranı ve I2C protokolüyle çalışan \textbf{ 0.96" OLED}  ekran üzerinden anlık olarak doğrulanmıştır. Oda sıcaklığında yapılan testlerde hedef sıcaklık eşiği yazılımsal olarak manipüle edilmiş, ayrıca buzlu su teması ile doğrudan fiziksel testler gerçekleştirilmiştir. Bu testler sonucunda \textbf{Bluetooth UART} komutunun terminale kayıpsız düştüğü ve bilgisayar üzerindeki Python kodunun müziği anında başlattığı doğrulanmıştır.


% ════════════════════════════════════════════════════════════════════════════
\section{Özellik 2: Temassız Jest Kontrolü ile Müzik Değiştirme}
% ════════════════════════════════════════════════════════════════════════════

\subsection{Çalışma Prensibi}

İki adet LDR (ışığa bağlı direnç) sensörü birbirinden $d = 5\,\mathrm{cm}$ mesafede
yerleştirilmiştir. El sensörün üzerinden geçtiğinde her LDR sırayla gölgelenir;
gölgelenme zaman farkı ($\Delta t$) elin hareket \textit{yönünü} belirler:

\begin{itemize}[leftmargin=*]
  \item $\Delta t > 0$ (Sol $\to$ Sağ): Bir sonraki şarkıya geç $\to$ \texttt{NEXT}
  \item $\Delta t < 0$ (Sağ $\to$ Sol): Önceki şarkıya geç $\to$ \texttt{PREV}
\end{itemize}

\subsection{Çift-Kapı Doğrulama Algoritması}

Ortamdaki ani ışık değişimlerinin yanlış jest olarak algılanmaması için
\textbf{çift-kapı doğrulama} mekanizması uygulanmıştır:

\begin{enumerate}[leftmargin=*, label=\arabic*.]
  \item İlk LDR tetiklendiğinde donanımsal zamanlayıcı başlatılır.
  \item İkinci LDR belirlenen $t_{\max} = 300\,\mathrm{ms}$ penceresi içinde
        tetiklenirse jest geçerli sayılır.
  \item Pencere süresi aşılırsa olay görmezden gelinir ve zamanlayıcı sıfırlanır.
\end{enumerate}

\subsection{Kritik Kod Parçacığı}

\begin{lstlisting}[style=cstyle, caption={Çift-kapı jest doğrulama mantığı}]
// GPIO ISR - LDR1 tetiklendi
void PORT3_IRQHandler(void) {
    uint32_t status = P3->IFG;
    P3->IFG = 0;

    if (status & LDR1_PIN) {
        t_ldr1 = Timer32_GetValue(TIMER32_0_BASE);
        ldr1_triggered = true;
    }
    if (status & LDR2_PIN && ldr1_triggered) {
        uint32_t delta = t_ldr1 - Timer32_GetValue(TIMER32_0_BASE);
        if (delta < MAX_GESTURE_TICKS) {
            // Gecerli jest - yonu belirle ve komut gonder
            if (ldr1_first) UART_SendString("NEXT\r\n");
            else            UART_SendString("PREV\r\n");
        }
        ldr1_triggered = false;
    }
}
\end{lstlisting}

\subsection{Test Sonuçları}

50 tekrarlı test gerçekleştirilmiş; doğruluk oranı \%94 olarak ölçülmüştür.
Hatalı algılamaların tamamı $\Delta t > 280\,\mathrm{ms}$ sınır bölgesindeki
yavaş el hareketlerinden kaynaklanmıştır.

% ════════════════════════════════════════════════════════════════════════════
\section{Özellik 3: Akustik Lamba Kontrolü}
% ════════════════════════════════════════════════════════════════════════════

\subsection{Çalışma Prensibi}

KY-037 mikrofon modülü analog çıkışı (A0) üzerinden Arduino Uno'ya sürekli ses verisi iletmektedir. Arduino her döngüde \texttt{analogRead()} ile 0--1023 aralığında bir değer okur ve bu değeri yazılımsal eşik (45) ile karşılaştırır. Eşiği aşan ses algılandığında alkış sayacı güncellenir; 1 saniyelik karar penceresi sonunda alkış sayısına göre işlem yapılır.

Tek alkışta lamba açılıp kapatılır. Çift alkışta ise lamba açıksa parlaklık kademesi değiştirilir. Parlaklık kontrolü Arduino'nun D9 PWM pini üzerinden IRF520 MOSFET modülüne sinyal gönderilerek sağlanır. MOSFET modülü 12V DC adaptörden beslenen LED lambayı sürer.

\subsection{Tek/Çift Alkış Algoritması}

\begin{enumerate}[leftmargin=*, label=\arabic*.]
  \item Ses eşiği aşıldığında \texttt{firstClap} zamanı kaydedilir, \texttt{clapCount} 1 yapılır.
  \item 1 saniyelik pencere içinde ikinci alkış gelirse \texttt{clapCount} 2 yapılır.
  \item 1 saniye dolunca karar verilir: 1 ise aç/kapat, 2 ise kademe değişimi.
  \item Ardışık darbeleri önlemek için 200 ms debounce uygulanır.
\end{enumerate}

\subsection{Parlaklık Kademeleri}

\begin{table}[H]
  \centering
  \caption{PWM Parlaklık Kademeleri}
  \begin{tabular}{ccc}
    \toprule
    \textbf{Kademe} & \textbf{PWM (0--255)} & \textbf{Parlaklık} \\
    \midrule
    1 & 26  & \%10  \\
    2 & 64  & \%25  \\
    3 & 128 & \%50  \\
    4 & 192 & \%75  \\
    5 & 255 & \%100 \\
    \bottomrule
  \end{tabular}
\end{table}

\subsection{Donanım Bağlantıları}

\begin{table}[H]
  \centering
  \caption{Akustik Lamba -- Pin Bağlantıları}
  \begin{tabular}{lll}
    \toprule
    \textbf{Bileşen} & \textbf{Arduino} & \textbf{Açıklama} \\
    \midrule
    KY-037 A0  & A0       & Analog ses okuma     \\
    KY-037 VCC & 5V       & Besleme              \\
    KY-037 GND & GND      & Toprak               \\
    IRF520 SIG & D9 (PWM) & PWM kontrol sinyali  \\
    IRF520 VCC & 5V       & Modül beslemesi      \\
    IRF520 GND & GND      & Ortak toprak         \\
    12V Adaptör & IRF520 V+/V-- & Lamba güç kaynağı \\
    LED Lamba  & IRF520 çıkışı & Anahtarlanan yük  \\
    \bottomrule
  \end{tabular}
\end{table}

\subsection{Kaynak Kodu}

\begin{lstlisting}[style=cstyle, caption={Akustik lamba kontrolü -- Arduino kodu}]
int soundPin = A0;
int lampPin = 9;
int level = 1;
int brightness[] = {0, 26, 64, 128, 192, 255};
unsigned long firstClap = 0;
unsigned long lastDetected = 0;
int clapCount = 0;
bool lampOn = false;

void setup() {
  pinMode(lampPin, OUTPUT);
  analogWrite(lampPin, 0);
  Serial.begin(9600);
}

void loop() {
  int value = analogRead(soundPin);
  unsigned long now = millis();

  if (value > 45 && now - lastDetected > 200) {
    lastDetected = now;
    if (clapCount == 0) {
      firstClap = now;
      clapCount = 1;
    } else {
      clapCount = 2;
    }
  }

  if (clapCount > 0 && now - firstClap > 1000) {
    if (clapCount == 1) {
      lampOn = !lampOn;
      analogWrite(lampPin, lampOn ? brightness[level] : 0);
    } 
    else if (clapCount == 2 && lampOn) {
      level++;
      if (level > 5) level = 1;
      analogWrite(lampPin, brightness[level]);
    }
    clapCount = 0;
  }
}
\end{lstlisting}

\subsection{Güvenlik Önlemleri}

220V AC şebeke yerine 12V DC adaptör kullanılarak elektrik çarpması riski ortadan kaldırılmıştır. Arduino 12V hattına doğrudan bağlı değildir; yük IRF520 üzerinden sürülmektedir.

% ════════════════════════════════════════════════════════════════════════════
\section{Donanım Bağlantı Şeması}
% ════════════════════════════════════════════════════════════════════════════

\begin{figure}[H]
  \centering
  % Gerçek devre şeması buraya eklenecek:
  % \includegraphics[width=0.9\textwidth]{devre.png}
  \fbox{\rule{0pt}{6cm}\rule{0.9\textwidth}{0pt}}
  \caption{Tam Devre Bağlantı Şeması}
  \label{fig:devre}
\end{figure}

\begin{table}[H]
  \centering
  \caption{MSP432 Pin Atama Tablosu}
  \label{tab:pinler}
  \begin{tabular}{llll}
    \toprule
    \textbf{Bileşen}  & \textbf{MSP432 Pini} & \textbf{Yön}  & \textbf{Protokol} \\
    \midrule
    DS18B20           & P2.0                 & Giriş/Çıkış   & 1-Wire       \\
    HC-05 TX          & P3.2 (UCA2RXD)       & Giriş         & UART         \\
    HC-05 RX          & P3.3 (UCA2TXD)       & Çıkış         & UART         \\
    LDR-1             & P3.4                 & Giriş         & GPIO INT     \\
    LDR-2             & P3.5                 & Giriş         & GPIO INT     \\
    KY-037 (DO)       & P5.1                 & Giriş         & GPIO INT     \\
    Röle              & P6.0                 & Çıkış         & GPIO         \\
    OLED SDA          & P6.4 (UCB1SDA)       & Giriş/Çıkış   & I2C          \\
    OLED SCL          & P6.5 (UCB1SCL)       & Çıkış         & I2C          \\
    \bottomrule
  \end{tabular}
\end{table}

% ════════════════════════════════════════════════════════════════════════════
\section{Test ve Doğrulama}
% ════════════════════════════════════════════════════════════════════════════

\subsection{Test Yöntemi}

Her özellik kendi alt sistemi içinde birim test edilmiş, ardından tüm sistem
entegre çalıştırılmıştır.

\subsection{Test Sonuçları Özeti}

\begin{table}[H]
  \centering
  \caption{Test Sonuçları}
  \begin{tabular}{lccc}
    \toprule
    \textbf{Test Senaryosu}          & \textbf{Deneme} & \textbf{Başarı} & \textbf{Oran} \\
    \midrule
    Sıcaklık eşiği algılama          & 20              & 20              & \%100 \\
    Bluetooth müzik başlatma          & 20              & 19              & \%95  \\
    Jest -- sağa kaydırma            & 25              & 24              & \%96  \\
    Jest -- sola kaydırma            & 25              & 23              & \%92  \\
    Lamba açma (el çırpma)           & 30              & 29              & \%97  \\
    Yanlış tetikleme (konuşma sesi)  & 20              & 0               & \%0   \\
    \bottomrule
  \end{tabular}
\end{table}

% ════════════════════════════════════════════════════════════════════════════
\section{Karşılaşılan Problemler ve Çözümler}
% ════════════════════════════════════════════════════════════════════════════

\begin{table}[H]
  \centering
  \caption{Karşılaşılan Teknik Sorunlar}
  \begin{tabular}{p{4cm}p{5.5cm}p{4cm}}
    \toprule
    \textbf{Sorun}           & \textbf{Belirtisi}                          & \textbf{Çözüm} \\
    \midrule
    Ortam gürültüsü & Konuşma sesi lambayı tetikliyordu & KY-037 hassasiyet ayarı + 200\,ms debounce \\\\
    Optik yanılsama          & Ani ışık değişimi jest olarak algılandı     & Çift-kapı doğrulama ($t_{\max}=300\,\mathrm{ms}$) \\
    UART iletişim gecikmesi  & Bluetooth geç bağlanıyordu                  & HC-05 baud hızı doğrulaması ve bekleme döngüsü \\
    Röle titremesi           & Kesme sınırında lamba sürekli toggling      & Yazılımsal kilit mekanizması \\
    \bottomrule
  \end{tabular}
\end{table}

% ════════════════════════════════════════════════════════════════════════════
\section{Zaman Planı ve Sprint Özeti}
% ════════════════════════════════════════════════════════════════════════════

\begin{table}[H]
  \centering
  \caption{Sprint Planı}
  \begin{tabular}{clp{8cm}}
    \toprule
    \textbf{Sprint} & \textbf{Tarih} & \textbf{Tamamlanan İşler} \\
    \midrule
    1 & Hafta 1 & Donanım temini, DS18B20 kalibrasyonu, LDR eşik testleri \\
    2 & Hafta 2 & Durum makinesi kodlaması, Timer kesmeleri, röle prototipleme \\
    3 & Hafta 3 & HC-05/macOS entegrasyonu, eş zamanlılık ve kararlılık testleri \\
    \bottomrule
  \end{tabular}
\end{table}

% ════════════════════════════════════════════════════════════════════════════
\section{Sonuç}
% ════════════════════════════════════════════════════════════════════════════

Bu proje kapsamında, MSP432 mikrodenetleyici üzerinde kesme tabanlı çoklu sensör
entegrasyonu başarıyla gerçeklenmiştir. Üç temel işlev (sıcaklık tabanlı müzik
başlatma, temassız jest kontrolü ve akustik lamba yönetimi) birbirinden bağımsız
alt sistemler olarak tasarlanmış; ortak bir durum makinesi çatısı altında
eş zamanlı ve kararlı biçimde çalıştırılmıştır. Test sonuçları, tüm özelliklerin
\%92'nin üzerinde başarı oranına ulaştığını göstermektedir.

% ════════════════════════════════════════════════════════════════════════════
\section*{Kaynaklar}
% ════════════════════════════════════════════════════════════════════════════

\begin{enumerate}[label={[\arabic*]}]
  \item Texas Instruments, \textit{MSP432P4xx SimpleLink™ Microcontrollers Technical Reference Manual}, 2017.
  \item Maxim Integrated, \textit{DS18B20 Programmable Resolution 1-Wire Digital Thermometer Datasheet}, 2019.
  \item HC-05 Bluetooth Module AT Command Reference, \url{https://components101.com/wireless/hc-05-bluetooth-module}.
  \item KY-037 Sound Sensor Module Datasheet, \url{https://arduinomodules.info/ky-037-microphone-sound-sensor-module/}.
\end{enumerate}

% ════════════════════════════════════════════════════════════════════════════
\appendix
\section{Tam Kaynak Kodu}
% ════════════════════════════════════════════════════════════════════════════

Projenin tüm kaynak kodları aşağıdaki GitHub deposunda yer almaktadır:\\
\url{https://github.com/KULLANICI-ADI/akilli-ekosistem}

% İsteğe bağlı: Kodları doğrudan buraya eklemek için aşağıdaki yapıyı kullanın:
% \lstinputlisting[style=cstyle, caption={main.c}]{../src/main.c}

\end{document}